%%%
%%% 17画
%%%
\section*{17画}\addcontentsline{toc}{section}{17画}\addcontentsline{loh}{figure}{\#\#\#\# 17画}

%%%%%%%%%% 徽 %%%%%%%%%%
\subsection*{徽}\addcontentsline{loh}{figure}{徽}

\begin{Entry}{徽}{17}{⼻}
  \begin{Phonetics}{徽}{hui1}
    \definition*{s.}{Huizhou (徽州), uma antiga prefeitura na província de Anhui (安徽)}
    \definition{adj.}{belo; esplêndido; magnífico}
    \definition[枚]{s.}{distintivo; emblema; insígnia; um símbolo ou sinal que representa um determinado grupo ou organização}
  \seealsoref{安徽}{an1hui1}
  \seealsoref{徽州}{hui1zhou1}
  \end{Phonetics}
\end{Entry}

\begin{Entry}{徽州}{17,6}{⼻,⼮}
  \begin{Phonetics}{徽州}{hui1zhou1}
    \definition*{s.}{Huizhou, antiga prefeitura que incluía o atual condado de Shexian (歙县) na província de Anhui (安徽), famosa como produtora de bastões de tinta}
  \seealsoref{安徽}{an1hui1}
  \seealsoref{歙县}{she4xian4}
  \end{Phonetics}
\end{Entry}

%%%%%%%%%% 戴 %%%%%%%%%%
\subsection*{戴}\addcontentsline{loh}{figure}{戴}

\begin{Entry}{戴}{17}{⼽}
  \begin{Phonetics}{戴}{dai4}[][HSK 4]
    \definition*{s.}{Sobrenome: Dai}
    \definition{v.}{usar/vestir (óculos, gravata, relógio de pulso, luvas); colocar objetos em sua cabeça, rosto, pescoço, peito, braços etc. | honrar; respeitar}
  \seealsoref{带}{dai4}
  \synonymref{承}{cheng2}
  \synonymref{顶}{ding3}
  \synonymref{护}{hu4}
  \synonymref{敬}{jing4}
  \synonymref{拥}{yong1}
  \synonymref{尊}{zun1}
  \antonymref{摘}{zhai1}
  \end{Phonetics}
\end{Entry}

%%%%%%%%%% 擦 %%%%%%%%%%
\subsection*{擦}\addcontentsline{loh}{figure}{擦}

\begin{Entry}{擦}{17}{⼿}
  \begin{Phonetics}{擦}{ca1}[][HSK 4]
    \definition{v.}{enxugar; esfregar; apagar; limpar; limpar esfregando com um pano, toalha de mão, etc. | espalhar sobre; colocar sobre | passar raspando | ralar (em pedaços); ralar frutas em um ralador para fazer fios finos}
  \seealsoref{擦拭}{ca1shi4}
  \synonymref{抹}{mo3}
  \synonymref{抹}{mo4}
  \synonymref{涂}{tu2}
  \end{Phonetics}
\end{Entry}

\begin{Entry}{擦拭}{17,9}{⼿,⼿}
  \begin{Phonetics}{擦拭}{ca1shi4}
    \definition{v.}{limpar; enxugar; esfregar; higienizar}
  \seealsoref{擦}{ca1}
  \end{Phonetics}
\end{Entry}

%%%%%%%%%% 曙 %%%%%%%%%%
\subsection*{曙}\addcontentsline{loh}{figure}{曙}

\begin{Entry}{曙}{17}{⽇}
  \begin{Phonetics}{曙}{shu3}
    \definition{s.}{Literário: amanhecer; alvorecer | Literário: o alvorecer de uma nova época (metáfora)}
  \end{Phonetics}
\end{Entry}

\begin{Entry}{曙光}{17,6}{⽇,⼉}
  \begin{Phonetics}{曙光}{shu3guang1}[][HSK 7-9]
    \definition{s.}{aurora; amanhecer; crepúsculo; primeira luz da manhã; a primeira luz do dia; crepúsculo matutino; luz do sol surgindo no horizonte | amanhecer; uma metáfora para um futuro brilhante que já está à vista}
  \end{Phonetics}
\end{Entry}

%%%%%%%%%% 朦 %%%%%%%%%%
\subsection*{朦}\addcontentsline{loh}{figure}{朦}

\begin{Entry}{朦}{17}{⽉}
  \begin{Phonetics}{朦}{meng2}
    \definition{adj.}{indistinto | pouco claro}
    \definition{v.}{enganar}
  \end{Phonetics}
\end{Entry}

\begin{Entry}{朦胧}{17,9}{⽉,⾁}
  \begin{Phonetics}{朦胧}{meng2long2}[][HSK 7-9]
    \definition{adj.}{fraco; nebuloso; obscuro; pouco claro; vago}
  \end{Phonetics}
\end{Entry}

%%%%%%%%%% 爵 %%%%%%%%%%
\subsection*{爵}\addcontentsline{loh}{figure}{爵}

\begin{Entry}{爵}{17}{⽖}
  \begin{Phonetics}{爵}{jue2}
    \definition{s.}{grau de nobreza; aristocracia | o título de nobreza | Arcaico: um antigo recipiente para vinho com três pés e uma alça em forma de laço}
  \end{Phonetics}
\end{Entry}

\begin{Entry}{爵士}{17,3}{⽖,⼠}
  \begin{Phonetics}{爵士}{jue2shi4}[][HSK 7-9]
    \definition*{s.}{\emph{Sir} (Senhor); Título honorífico para cavalheiros da Ordem do Império Britânico}
    \definition{s.}{cavaleiro; o título mais baixo de uma monarquia europeia | Empréstimo linguístico: jazz}[它赞助大型爵士音乐会。===A empresa patrocina concertos de jazz de grande escala.]
  \end{Phonetics}
\end{Entry}

%%%%%%%%%% 癌 %%%%%%%%%%
\subsection*{癌}\addcontentsline{loh}{figure}{癌}

\begin{Entry}{癌}{17}{⽧}
  \begin{Phonetics}{癌}{ai2}[][HSK 7-9]
    \definition{s.}{câncer; carcinoma; tumor maligno}
  \end{Phonetics}
\end{Entry}

\begin{Entry}{癌症}{17,10}{⽧,⽧}
  \begin{Phonetics}{癌症}{ai2zheng4}[][HSK 7-9]
    \definition[种]{s.}{câncer; tumores malignos no corpo}
  \end{Phonetics}
\end{Entry}

%%%%%%%%%% 瞧 %%%%%%%%%%
\subsection*{瞧}\addcontentsline{loh}{figure}{瞧}

\begin{Entry}{瞧}{17}{⽬}
  \begin{Phonetics}{瞧}{qiao2}[][HSK 5]
    \definition{v.}{ver; olhar | tratar; diagnosticar e tratar | ver; visitar; fazer uma visita}
  \synonymref{瞅}{chou3}
  \synonymref{顾}{gu4}
  \synonymref{观}{guan1}
  \synonymref{看}{kan4}
  \synonymref{视}{shi4}
  \synonymref{望}{wang4}
  \end{Phonetics}
\end{Entry}

\begin{Entry}{瞧不起}{17,4,10}{⽬,⼀,⾛}
  \begin{Phonetics}{瞧不起}{qiao2bu5qi3}[][HSK 7-9]
    \definition{v.}{olhar com desdém para alguém; menosprezar alguém; torcer o nariz para alguém; desprezar}
  \end{Phonetics}
\end{Entry}

%%%%%%%%%% 瞩 %%%%%%%%%%
\subsection*{瞩}\addcontentsline{loh}{figure}{瞩}

\begin{Entry}{瞩}{17}{⽬}
  \begin{Phonetics}{瞩}{zhu3}
    \definition{v.}{olhar fixamente; fitar o olhar | concentrar o olhar em}
  \end{Phonetics}
\end{Entry}

\begin{Entry}{瞩目}{17,5}{⽬,⽬}
  \begin{Phonetics}{瞩目}{zhu3mu4}[][HSK 7-9]
    \definition{v.}{fixar o olhar em; concentrar a atenção em; contemplar; prestar atenção}
  \synonymref{注视}{zhu4shi4}
  \synonymref{注意}{zhu4yi4}
  \end{Phonetics}
\end{Entry}

%%%%%%%%%% 瞪 %%%%%%%%%%
\subsection*{瞪}\addcontentsline{loh}{figure}{瞪}

\begin{Entry}{瞪}{17}{⽬}
  \begin{Phonetics}{瞪}{deng4}[][HSK 7-9]
    \definition{v.}{abrir bem os olhos; encarar; olhar fixamente com os olhos bem abertos; expressar insatisfação}
  \end{Phonetics}
\end{Entry}

%%%%%%%%%% 瞬 %%%%%%%%%%
\subsection*{瞬}\addcontentsline{loh}{figure}{瞬}

\begin{Entry}{瞬}{17}{⽬}
  \begin{Phonetics}{瞬}{shun4}
    \definition{s.}{piscadela; cintilação; num piscar de olhos}
    \definition{v.}{piscar; cintilar}
  \end{Phonetics}
\end{Entry}

\begin{Entry}{瞬间}{17,7}{⽬,⾨}
  \begin{Phonetics}{瞬间}{shun4jian1}[][HSK 7-9]
    \definition{s.}{momento; instante; minuto; num piscar de olhos; num abrir e fechar de olhos; em um tempo extremamente curto}
  \antonymref{长期}{chang2qi1}
  \antonymref{漫长}{man4chang2}
  \antonymref{永恒}{yong3heng2}
  \antonymref{永远}{yong3yuan3}
  \end{Phonetics}
\end{Entry}

%%%%%%%%%% 礁 %%%%%%%%%%
\subsection*{礁}\addcontentsline{loh}{figure}{礁}

\begin{Entry}{礁}{17}{⽯}
  \begin{Phonetics}{礁}{jiao1}
    \definition{s.}{recife; rocha}
  \end{Phonetics}
\end{Entry}

\begin{Entry}{礁石}{17,5}{⽯,⽯}
  \begin{Phonetics}{礁石}{jiao1shi2}[][HSK 7-9]
    \definition{s.}{recife; rocha}
  \end{Phonetics}
\end{Entry}

%%%%%%%%%% 窾 %%%%%%%%%%
\subsection*{窾}\addcontentsline{loh}{figure}{窾}

\begin{Entry}{窾}{17}{⽳}
  \begin{Phonetics}{窾}{cuan4}
    \definition{adj.}{vazio | seco | destituído; pobre}
    \definition{s.}{buraco | lei}
    \definition{v.}{esconder}
  \end{Phonetics}
  \begin{Phonetics}{窾}{kuan3}
    \definition{adj.}{oco}
    \definition{s.}{rachadura; cavidade | (onomatopéia) água batendo na rocha}
    \definition{v.}{escavar um buraco}
  \end{Phonetics}
\end{Entry}

%%%%%%%%%% 簇 %%%%%%%%%%
\subsection*{簇}\addcontentsline{loh}{figure}{簇}

\begin{Entry}{簇}{17}{⽵}
  \begin{Phonetics}{簇}{cu4}
    \definition{clas.}{aglomerado; grupo; usado para pessoas ou coisas que se reúnem em grupos ou pilhas}
    \definition{s.}{pilha; aglomerado; buquê}
    \definition{v.}{aglomerar"-se; formar um aglomerado; empilhar}
  \end{Phonetics}
\end{Entry}

\begin{Entry}{簇拥}{17,8}{⽵,⼿}
  \begin{Phonetics}{簇拥}{cu4yong1}[][HSK 7-9]
    \definition{v.}{aglomerar"-se em volta de  | escoltar}
  \end{Phonetics}
\end{Entry}

%%%%%%%%%% 糟 %%%%%%%%%%
\subsection*{糟}\addcontentsline{loh}{figure}{糟}

\begin{Entry}{糟}{17}{⽶}
  \begin{Phonetics}{糟}{zao1}[][HSK 5]
    \definition{adj.}{pobre; apodrecido; deteriorado | estragado; em uma bagunça; em um estado miserável (terrível) | (situação ou circunstância) ruim; desfavorável}
    \definition{s.}{resíduos de destilação de bebidas alcoólicas; resíduos do processo de fermentação do vinho}
    \definition{v.}{marinar alimentos em vinho ou mosto | desperdiçar; estragar; destruir}
  \synonymref{差}{cha4}
  \synonymref{费}{fei4}
  \synonymref{腐}{fu3}
  \synonymref{坏}{huai4}
  \synonymref{烂}{lan4}
  \synonymref{损}{sun3}
  \antonymref{好}{hao3}
  \end{Phonetics}
\end{Entry}

\begin{Entry}{糟糕}{17,16}{⽶,⽶}
  \begin{Phonetics}{糟糕}{zao1gao1}[][HSK 5]
    \definition{adj.}{(corpo, situação, etc.) muito ruim, péssimo}
    \definition{interj.}{``Que terrível!''; ``Que má sorte!''; ``Muito ruim!''}
  \seealsoref{倒霉}{dao3//mei2}
  \synonymref{恶劣}{e4lie4}
  \antonymref{出色}{chu1se4}
  \antonymref{精彩}{jing1cai3}
  \antonymref{美好}{mei3hao3}
  \antonymref{漂亮}{piao4liang5}
  \end{Phonetics}
\end{Entry}

%%%%%%%%%% 繁 %%%%%%%%%%
\subsection*{繁}\addcontentsline{loh}{figure}{繁}

\begin{Entry}{繁}{17}{⽷}
  \begin{Phonetics}{繁}{fan2}
    \definition{adj.}{em grande número; numerosos; múltiplos | em grande número; numerosos; complexos; complicado}
    \definition{v.}{propagar; multiplicar}
  \antonymref{简}{jian3}
  \end{Phonetics}
\end{Entry}

\begin{Entry}{繁华}{17,6}{⽷,⼗}
  \begin{Phonetics}{繁华}{fan2hua2}[][HSK 7-9]
    \definition{adj.}{ocupado; agitado; próspero; florescente; (cidade, mercado de rua) movimentado e próspero}
  \end{Phonetics}
\end{Entry}

\begin{Entry}{繁忙}{17,6}{⽷,⼼}
  \begin{Phonetics}{繁忙}{fan2mang2}[][HSK 7-9]
    \definition{adj.}{ocupado; agitado; muita coisa para fazer, pouco tempo livre}
  \end{Phonetics}
\end{Entry}

\begin{Entry}{繁体字}{17,7,6}{⽷,⼈,⼦}
  \begin{Phonetics}{繁体字}{fan2ti3zi4}[][HSK 7-9]
    \definition{s.}{forma complexa tradicional de um caractere chinês simplificado; caracteres chineses com mais traços que foram substituídos por caracteres simplificados}
  \antonymref{简体字}{jian3ti3zi4}
  \end{Phonetics}
\end{Entry}

\begin{Entry}{繁花}{17,7}{⽷,⾋}
  \begin{Phonetics}{繁花}{fan2hua1}
    \definition{s.}{Literário: flores desabrochadas; flores de cores diferentes; flores exuberantes; uma massa de flores}
  \end{Phonetics}
\end{Entry}

\begin{Entry}{繁荣}{17,9}{⽷,⾋}
  \begin{Phonetics}{繁荣}{fan2rong2}[][HSK 5]
    \definition{adj.}{florescente; próspero}
    \definition{v.}{promover; prosperar}
  \seealsoref{繁华}{fan2hua2}
  \synonymref{发达}{fa1da2}
  \synonymref{发展}{fa1zhan3}
  \synonymref{茂盛}{mao4sheng4}
  \synonymref{蓬勃}{peng2bo2}
  \synonymref{热闹}{re4nao5}
  \synonymref{旺盛}{wang4sheng4}
  \synonymref{兴旺}{xing1wang4}
  \antonymref{废墟}{fei4xu1}
  \antonymref{荒凉}{huang1liang2}
  \antonymref{凄凉}{qi1liang2}
  \antonymref{萧条}{xiao1tiao2}
  \end{Phonetics}
\end{Entry}

\begin{Entry}{繁重}{17,9}{⽷,⾥}
  \begin{Phonetics}{繁重}{fan2zhong4}[][HSK 7-9]
    \definition{adj.}{pesado; oneroso; árduo; penoso; (trabalho, tarefas) muitas e pesadas}[我每天都有繁重的工作。===Tenho uma carga de trabalho pesada todos os dias.]
  \end{Phonetics}
\end{Entry}

\begin{Entry}{繁殖}{17,12}{⽷,⽍}
  \begin{Phonetics}{繁殖}{fan2zhi2}[][HSK 6]
    \definition{v.}{criar; reproduzir; propagar; multiplicar; os organismos produzem novos indivíduos}
  \synonymref{孵化}{fu1hua4}
  \synonymref{衍生}{yan3sheng1}
  \antonymref{灭绝}{mie4jue2}
  \end{Phonetics}
\end{Entry}

%%%%%%%%%% 藏 %%%%%%%%%%
\subsection*{藏}\addcontentsline{loh}{figure}{藏}

\begin{Entry}{藏}{17}{⾋}
  \begin{Phonetics}{藏}{cang2}[][HSK 6]
    \definition*{s.}{Sobrenome: Cang}
    \definition{v.}{esconder; ocultar; esconder da vista | armazenar; coletar; colocar de lado}
  \seealsoref{躲}{duo3}
  \synonymref{含}{han2}
  \synonymref{匿}{ni4}
  \synonymref{隐}{yin3}
  \antonymref{露}{lou4}
  \antonymref{露}{lu4}
  \antonymref{显}{xian3}
  \end{Phonetics}
  \begin{Phonetics}{藏}{zang4}
    \definition*{s.}{Escrituras budistas ou taoístas; um termo geral para clássicos budistas ou taoístas | Região Autônoma do Tibete, 西藏}
    \definition{s.}{depósito; local de armazenamento; armazém; local onde grandes quantidades de coisas são armazenadas | os tibetanos, 藏族; grupo étnico Zang (ou tibetano)}
  \seealsoref{西藏}{xi1zang4}
  \seealsoref{藏族}{zang4zu2}
  \end{Phonetics}
\end{Entry}

\begin{Entry}{藏身}{17,7}{⾋,⾝}
  \begin{Phonetics}{藏身}{cang2shen1}[][HSK 7-9]
    \definition{v.}{esconder"-se; esconder | abrigar"-se; estabelecer"-se}
  \end{Phonetics}
\end{Entry}

\begin{Entry}{藏品}{17,9}{⾋,⼝}
  \begin{Phonetics}{藏品}{cang2pin3}[][HSK 7-9]
    \definition[件]{s.}{artigos coletados; coleção | item de colecionador | peça de museu | objeto precioso}
  \end{Phonetics}
\end{Entry}

\begin{Entry}{藏匿}{17,10}{⾋,⼖}
  \begin{Phonetics}{藏匿}{cang2ni4}[][HSK 7-9]
    \definition{v.}{esconder; esconder"-se | abrigar}
  \end{Phonetics}
\end{Entry}

\begin{Entry}{藏族}{17,11}{⾋,⽅}
  \begin{Phonetics}{藏族}{zang4zu2}
    \definition*{s.}{Etnia Zang (ou tibetana); Os Zangs (ou tibetanos) , distribuídos pela Região Autônoma do Tibete e pelas províncias de Qinghai, Sichuan, Gansu e Yunnan}
  \end{Phonetics}
\end{Entry}

%%%%%%%%%% 螺 %%%%%%%%%%
\subsection*{螺}\addcontentsline{loh}{figure}{螺}

\begin{Entry}{螺}{17}{⾍}
  \begin{Phonetics}{螺}{luo2}
    \definition{s.}{concha espiral; caracol | espiral (na impressão digital) | concha (moluscos)}
  \end{Phonetics}
\end{Entry}

\begin{Entry}{螺丝}{17,5}{⾍,⼀}
  \begin{Phonetics}{螺丝}{luo2si1}[][HSK 7-9]
    \definition[个,颗]{s.}{parafuso | metáfora para pessoa ou coisa comum, porém indispensável; peças que desempenham um papel importante na organização}
  \end{Phonetics}
\end{Entry}

%%%%%%%%%% 豁 %%%%%%%%%%
\subsection*{豁}\addcontentsline{loh}{figure}{豁}

\begin{Entry}{豁}{17}{⾕}
  \begin{Phonetics}{豁}{hua2}
    \definition{v.}{jogar o jogo de adivinhação de dedos chinês (Huoquan); o mesmo que 划拳}[豁拳规则很简单。===As regras do Huoquan são muito simples.]
  \seealsoref{划拳}{hua2quan2}
  \end{Phonetics}
  \begin{Phonetics}{豁}{huo1}[][HSK 7-9]
    \definition{v.}{cortar; quebrar; rachar; dividir | sacrificar; desistir; pagar o preço cruelmente}
  \end{Phonetics}
  \begin{Phonetics}{豁}{huo4}
    \definition{adj.}{claro; aberto; de mente aberta; generoso}
    \definition{v.}{isentar; remeter}
  \end{Phonetics}
\end{Entry}

\begin{Entry}{豁出去}{17,5,5}{⾕,⼐,⼛}
  \begin{Phonetics}{豁出去}{huo1//chu5qu4}[][HSK 7-9]
    \definition{v.+compl.}{seguir em frente independentemente; pronto para arriscar tudo}[她豁出去所有，去追逐梦想。===Ela arriscou tudo para perseguir seu sonho.]
  \end{Phonetics}
\end{Entry}

\begin{Entry}{豁达}{17,6}{⾕,⾡}
  \begin{Phonetics}{豁达}{huo4da2}[][HSK 7-9]
    \definition{adj.}{otimista; de mente aberta; aberto e claro}
  \end{Phonetics}
\end{Entry}

%%%%%%%%%% 赡 %%%%%%%%%%
\subsection*{赡}\addcontentsline{loh}{figure}{赡}

\begin{Entry}{赡}{17}{⾙}
  \begin{Phonetics}{赡}{shan4}[][HSK 7-9]
    \definition*{s.}{Sobrenome: Shan}
    \definition{adj.}{Literário: suficiente; abundante; adequado | Literário: refinado, rico}
    \definition{v.}{apoiar; prover para}
  \end{Phonetics}
\end{Entry}

\begin{Entry}{赡养}{17,9}{⾙,⼋}
  \begin{Phonetics}{赡养}{shan4yang3}[][HSK 7-9]
    \definition{v.}{apoiar; prover para; prover as necessidades básicas refere"-se especificamente à ajuda que os filhos dão aos pais, tanto materialmente quanto no dia a dia}
  \synonymref{抚养}{fu3yang3}
  \end{Phonetics}
\end{Entry}

%%%%%%%%%% 赢 %%%%%%%%%%
\subsection*{赢}\addcontentsline{loh}{figure}{赢}

\begin{Entry}{赢}{17}{⾙}
  \begin{Phonetics}{赢}{ying2}[][HSK 3]
    \definition{v.}{vencer; derrotar | ganhar; lucrar}
  \seealsoref{胜}{sheng4}
  \synonymref{券}{quan4}
  \synonymref{赚}{zhuan4}
  \antonymref{败}{bai4}
  \antonymref{负}{fu4}
  \antonymref{赔}{pei2}
  \antonymref{输}{shu1}
  \end{Phonetics}
\end{Entry}

\begin{Entry}{赢家}{17,10}{⾙,⼧}
  \begin{Phonetics}{赢家}{ying2jia1}[][HSK 7-9]
    \definition[个]{s.}{vencedor; o vencedor de uma aposta ou competição}
  \antonymref{输家}{shu1jia1}
  \end{Phonetics}
\end{Entry}

\begin{Entry}{赢得}{17,11}{⾙,⼻}
  \begin{Phonetics}{赢得}{ying2de2}[][HSK 4]
    \definition{v.}{ganhar; obter; conquistar; assegurar; garantir}
  \synonymref{得到}{de2dao4}
  \synonymref{获得}{huo4de2}
  \synonymref{取得}{qu3de2}
  \antonymref{失掉}{shi1diao4}
  \antonymref{失去}{shi1qu4}
  \end{Phonetics}
\end{Entry}

%%%%%%%%%% 辫 %%%%%%%%%%
\subsection*{辫}\addcontentsline{loh}{figure}{辫}

\begin{Entry}{辫}{17}{⾟}
  \begin{Phonetics}{辫}{bian4}
    \definition{s.}{trança; rabo de cavalo | para coisas como uma trança}
  \end{Phonetics}
\end{Entry}

\begin{Entry}{辫子}{17,3}{⾟,⼦}
  \begin{Phonetics}{辫子}{bian4zi5}[][HSK 7-9]
    \definition[条,根,种]{s.}{trança; rabo de cavalo; uma mecha de cabelo presa reta ou trançada em seções | Metafórico: erro; deficiência; fraqueza | coisas parecidas com tranças}
  \end{Phonetics}
\end{Entry}

%%%%%%%%%% 邉 %%%%%%%%%%
\subsection*{邉}\addcontentsline{loh}{figure}{邉}

\begin{Entry}{邉}{17}{⾡}
  \begin{Phonetics}{邉}{bian1}
    \variantof{边}
  \end{Phonetics}
\end{Entry}

%%%%%%%%%% 镫 %%%%%%%%%%
\subsection*{镫}\addcontentsline{loh}{figure}{镫}

\begin{Entry}{镫}{17}{⾦}
  \begin{Phonetics}{镫}{deng4}
    \definition{s.}{estribo; os apoios para os pés que ficam pendurados em ambos os lados do selim são, em sua maioria, feitos de ferro}
  \end{Phonetics}
\end{Entry}

\begin{Entry}{镫骨}{17,9}{⾦,⾻}
  \begin{Phonetics}{镫骨}{deng4gu3}
    \definition{s.}{estribo, osso do ouvido médio que transmite a vibração sonora para o ouvido interno; pedistíbulo; um dos ossículos auditivos, com formato semelhante a um estribo, conecta"-se à bigorna externamente e ao ouvido interno internamente}
  \end{Phonetics}
\end{Entry}

%%%%%%%%%% 霜 %%%%%%%%%%
\subsection*{霜}\addcontentsline{loh}{figure}{霜}

\begin{Entry}{霜}{17}{⾬}
  \begin{Phonetics}{霜}{shuang1}[][HSK 7-9]
    \definition*{s.}{Sobrenome: Shuang}
    \definition{adj.}{branco; grisalho}
    \definition{s.}{geada | pó branco ou creme espalhado por uma superfície | glacê | creme para os olhos}
  \end{Phonetics}
\end{Entry}

%%%%%%%%%% 鞠 %%%%%%%%%%
\subsection*{鞠}\addcontentsline{loh}{figure}{鞠}

\begin{Entry}{鞠}{17}{⾰}
  \begin{Phonetics}{鞠}{ju1}
    \definition*{s.}{Sobrenome: Ju}
    \definition{s.}{arco | Arcaico: uma bola de futebol; um tipo antigo de bola; bola usada para jogar}
    \definition{v.}{Literátio: criar; educar; nutrir | obrar; curvar}
  \end{Phonetics}
\end{Entry}

\begin{Entry}{鞠躬}{17,10}{⾰,⾝}
  \begin{Phonetics}{鞠躬}{ju1//gong1}[][HSK 7-9]
    \definition{v.+compl.}{curvar"-se; inclinar"-se}
  \end{Phonetics}
\end{Entry}

%%%%%%%%%% 骤 %%%%%%%%%%
\subsection*{骤}\addcontentsline{loh}{figure}{骤}

\begin{Entry}{骤}{17}{⾺}
  \begin{Phonetics}{骤}{zhou4}
    \definition{adj.}{súbito; abrupto | rápido; veloz}
    \definition{adv.}{subitamente; abruptamente}
    \definition{v.}{(um cavalo) trotar; andar rápido}
  \end{Phonetics}
\end{Entry}

\begin{Entry}{骤然}{17,12}{⾺,⽕}
  \begin{Phonetics}{骤然}{zhou4ran2}[][HSK 7-9]
    \definition{adv.}{Literário: de repente; abruptamente; inesperadamente; de uma só vez}
  \synonymref{忽然}{hu1ran2}
  \synonymref{猛然}{meng3ran2}
  \synonymref{突然}{tu1ran2}
  \antonymref{缓慢}{huan3man4}
  \antonymref{渐渐}{jian4jian4}
  \end{Phonetics}
\end{Entry}

%%%%%%%%%% 鳄 %%%%%%%%%%
\subsection*{鳄}\addcontentsline{loh}{figure}{鳄}

\begin{Entry}{鳄}{17}{⿂}
  \begin{Phonetics}{鳄}{e4}
    \definition{s.}{crocodilo;  jacaré}
  \end{Phonetics}
\end{Entry}

\begin{Entry}{鳄鱼}{17,8}{⿂,⿂}
  \begin{Phonetics}{鳄鱼}{e4yu2}[][HSK 7-9]
    \definition[只,群,些]{s.}{crocodilo; jacaré}
  \end{Phonetics}
\end{Entry}

%%%%%%%%%% 黏 %%%%%%%%%%
\subsection*{黏}\addcontentsline{loh}{figure}{黏}

\begin{Entry}{黏}{17}{⿉}
  \begin{Phonetics}{黏}{nian2}[][HSK 7-9]
    \definition{adj.}{pegajoso; adesivo; glutinoso}
  \end{Phonetics}
\end{Entry}

%%%%%%%%%% 龠 %%%%%%%%%%
\subsection*{龠}\addcontentsline{loh}{figure}{龠}

\begin{Entry}{龠}{17}{⿕}[Kangxi 214]
  \begin{Phonetics}{龠}{yue4}
    \definition{clas.}{yue, uma unidade de medida seca para grãos (= 0,5 decilitro)}
    \definition{s.}{uma flauta curta antiga}
  \end{Phonetics}
\end{Entry}

%%%%% EOF %%%%%

